\renewcommand{\baselinestretch}{1.05}\normalfont
\chapter*{Prefacio}
\thispagestyle{empty}
\vfill

El presente texto es una edición recompuesta
\ifthenelse{\boolean{orig}}{\ignorespaces}{y anotada}
del libro «Recubrimientos étales y grupo fundamental», Lecture Notes in Mathematics, 224, Springer-Verlag, Berlín-Heidelberg-Nueva York, 1971, de Alexander Grothendieck y otros.

La composición en \LaTeX 2e fue realizada por voluntarios que participaron en un proyecto dirigido por Bas Edixhoven, sobre el cual pueden encontrarse más detalles en \url{http://www.math.leidenuniv.nl/~edix/}. La maquetación fue terminada por la Société mathématique de France.

\ifthenelse{\boolean{orig}}
{Se supone que esta versión reproduce el \emph{texto original}. Existe otra versión en la que se han corregido algunos errores tipográficos y Michel Raynaud ha realizado algunas actualizaciones. Esta última es publicada por la Société mathématique de France (\url{http://smf.emath.fr/}) en la serie «Documents Mathématiques».}
{Esta es una \emph{versión ligeramente corregida} del texto original. Es publicada por la Société mathématique de France (\url{http://smf.emath.fr/}) en la serie «Documents Mathématiques». Michel Raynaud ha añadido algunas observaciones de actualización. Estas están delimitadas por corchetes \lcrochetbf\,\rcrochetbf\ e indicadas con el símbolo (MR) (observaciones en las páginas \pageref{X.2.14}, \pageref{XI.1.4}, \pageref{XII.5.6}, \pageref{XIII.2.13} y nota \Ref{III.6.6.p24} en la página \pageref{III.6.6.p24}). Para uniformar la notación, el cuerpo residual de un punto $x$ se denota por $\kres(x)$, y el de un anillo local $A$, por $\kres(A)$.\par
Existe asimismo una versión electrónica que pretende reproducir el texto original.}

Las dos versiones se producen a partir de un único archivo fuente, alojado en el servidor de prepublicaciones de arXiv.org, \url{http://arxiv.org/}. Todas las diferencias entre ambas versiones están documentadas en el archivo fuente.

La antigua numeración de páginas se incorpora en el margen (el número~$n$ indica el comienzo de la página~$n$).

\vfill
\newpage

\begin{otherlanguage}{english}
The present text is a new
\ifthenelse{\boolean{orig}}{\ignorespaces}{updated}
edition of the book ``Rev\^etements \'Etales et Groupe Fondamental'', Lecture Notes in Mathematics, 224, Springer-Verlag, Berlin-Heidelberg-New York, 1971, by Alexander Grothendieck et al.

Typesetting in \LaTeX 2e was done by volunteers, participating in a project directed by Bas Edixhoven, on which one can find more details at \url{http://www.math.leidenuniv.nl/~edix/}. It has been completed by the Société mathématique de France.

\ifthenelse{\boolean{orig}}
{This version is meant to reproduce the \emph{original text}. There is another version, in which some typograpical mistakes have been corrected and some notes have been added by Michel Raynaud. The latter has been published by the Société mathématique de France (\url{http://smf.emath.fr/}) as a volume of the series ``Documents Mathématiques''.}
{This is a slightly \emph{corrected version} of the original text. It is published by the Société mathématique de France (\url{http://smf.emath.fr/}) as a volume of the series ``Documents Mathématiques''. Some updating remarks have been added by Michel Raynaud. These are bounded by brackets \lcrochetbf\,\rcrochetbf\ and indicated by the symbol (MR) (Remarks on pages \pageref{X.2.14}, \pageref{XI.1.4}, \pageref{XII.5.6}, \pageref{XIII.2.13} and footnote \Ref{III.6.6.p24} on page \pageref{III.6.6.p24}). In order to unify notation, the residue field of a point $x$ is denoted by $\kres(x)$ and the residue field of a local ring $A$ is denoted by $\kres(A)$.\par
There also exists an electronic version that is meant to reproduce the original text.}

Both versions are produced from the same source file, that can be downloaded from the arXiv.org e-Print server at \url{http://arxiv.org/}. All differences between the two versions are documented in the source file.

The old page numbering is incorporated in the margin (the number~$n$ marking the beginning of page~$n$).
\end{otherlanguage}

\vfill
\cleardoublepage
